Considera la funció
\[
h(x)=\begin{cases}
  2x+1 & \si{x<0},\\[3pt]
  e^{x} & \si{0\le x\le 2},\\[3pt]
  \dfrac{x-2}{x^2-5x+6} & \si{x>2}.
\end{cases}
\]

\begin{apartats}

\begin{nomesllarg}
\apartat{0,75}
Estudia la continuïtat de $h$ en $x=0$.

\begin{solucio}
$\lim_{x\to0^-}h(x)=1$, $\lim_{x\to0^+}h(x)=e^0=1$ i $h(0)=1$. \textbf{És contínua.}
\end{solucio}
\end{nomesllarg}

\apartat[1,25]{1}
Estudia la continuïtat de $h$ en $x=2$ i classifica la discontinuïtat, si n'hi ha.

\begin{solucio}
Per a $x>2$, $\dfrac{x-2}{(x-2)(x-3)}=\dfrac{1}{x-3}$.\\
$h(2)=e^2$ i $\lim_{x\to2^-}h(x)=e^2$, però $\lim_{x\to2^+}h(x)=\dfrac{1}{2-3}=-1$.
Laterals finits i diferents: \textbf{salt finit}.
\end{solucio}

\begin{tria}{trosos-tasca-c}
\itemtria{original}{0,75}{1,25}
Estudia la continuïtat de $h$ en $x=3$ i classifica la discontinuïtat, si n'hi ha.

\begin{solucio}
$h(3)$ no existeix. Per a $x>2$, $h(x)=\dfrac{1}{x-3}$, i els laterals valen
$-\infty$ (per l'esquerra) i $+\infty$ (per la dreta): \textbf{salt infinit} (asímptota
vertical $x=3$).
\end{solucio}

\itemtria{alternativa}{0,75}{1,25}
Calcula, si existeixen, aquests valors:
\[
h(-1) \qquad h(0) \qquad h(1) \qquad h(3)
\]

\begin{solucio}
$h(-1)=2(-1)+1=-1$ (primera branca). $h(0)=e^0=1$ i $h(1)=e^1=e$ (segona branca).
$h(3)$ \textbf{no existeix}: la tercera branca val $\dfrac{x-2}{(x-2)(x-3)}=\dfrac1{x-3}$,
que no es defineix en $x=3$.
\end{solucio}
\end{tria}

\end{apartats}
